\predpoglavje{Zahvala}

Pisanje zahvale ali posvetila je neobvezno. Pri zahvali napišite zahvalo tistim, ki so pomagali pri nastanku dela in brez katerih delo ne bi nastalo v takšni obliki, kot je. Navadno se je potrebno zahvaliti v prvi vrsti mentorju, somentorju in inštituciji, ki je morda finančno ali kako drugače podprla izvedbo doktorskega dela. Nato sledijo asistenti, ostali sodelavci in vaši kolegi, ki so pomagali pri teoretičnem in eksperimentalnem delu. Na koncu se po navadi zahvalimo še družini.
